%-------------------------
% Resume in Latex
% Author : Jake Gutierrez
% Based off: https://github.com/sb2nov/resume
% License : MIT
%------------------------

\documentclass[letterpaper,11pt]{article}

\usepackage{latexsym}
\usepackage[empty]{fullpage}
\usepackage{titlesec}
\usepackage{marvosym}
\usepackage[usenames,dvipsnames]{color}
\usepackage{verbatim}
\usepackage{enumitem}
\usepackage[hidelinks]{hyperref}
\usepackage{fancyhdr}
\usepackage[english]{babel}
\usepackage{tabularx}
\usepackage{iftex}
\ifPDFTeX
  \input{glyphtounicode}
  \pdfgentounicode=1
\fi


%----------FONT OPTIONS----------
% sans-serif
% \usepackage[sfdefault]{FiraSans}
% \usepackage[sfdefault]{roboto}
% \usepackage[sfdefault]{noto-sans}
% \usepackage[default]{sourcesanspro}

% serif
% \usepackage{CormorantGaramond}
% \usepackage{charter}


\pagestyle{fancy}
\fancyhf{} % clear all header and footer fields
\fancyfoot{}
\renewcommand{\headrulewidth}{0pt}
\renewcommand{\footrulewidth}{0pt}

% Adjust margins
\addtolength{\oddsidemargin}{-0.5in}
\addtolength{\evensidemargin}{-0.5in}
\addtolength{\textwidth}{1in}
\addtolength{\topmargin}{-.7in}
\addtolength{\textheight}{1.4in}

\urlstyle{same}

\raggedbottom
\raggedright
\setlength{\tabcolsep}{0in}

% Sections formatting
\titleformat{\section}{
  \vspace{-4pt}\scshape\raggedright\large\bfseries
}{}{0em}{}[\color{black}\titlerule \vspace{-5pt}]


%-------------------------
% Custom commands
\newcommand{\resumeItem}[1]{
  \item
    {#1 \vspace{-2pt}}
}

\newcommand{\resumeSubheading}[4]{
  \vspace{-2pt}\item
    \begin{tabular*}{0.97\textwidth}[t]{l@{\extracolsep{\fill}}r}
      \textbf{#1} & #2 \\
      \textit{#3} & {#4} \\
    \end{tabular*}\vspace{-7pt}
}

\newcommand{\resumeSubSubheading}[2]{
    \item
    \begin{tabular*}{0.97\textwidth}{l@{\extracolsep{\fill}}r}
      \textit{#1} & \textit{#2} \\
    \end{tabular*}\vspace{-7pt}
}

\newcommand{\resumeProjectHeading}[2]{
    \item
    \begin{tabular*}{0.97\textwidth}{l@{\extracolsep{\fill}}r}
      #1 & #2 \\
    \end{tabular*}\vspace{-7pt}
}

\newcommand{\resumeSubItem}[1]{\resumeItem{#1}\vspace{-4pt}}

\renewcommand\labelitemii{$\vcenter{\hbox{\tiny$\bullet$}}$}

\newcommand{\resumeSubHeadingListStart}{\begin{itemize}[leftmargin=0.15in, label={}]}
\newcommand{\resumeSubHeadingListEnd}{\end{itemize}}
\newcommand{\resumeItemListStart}{\begin{itemize}[topsep=2pt, itemsep=2pt, parsep=0pt]}
\newcommand{\resumeItemListEnd}{\end{itemize}\vspace{-5pt}}

%-------------------------------------------
%%%%%%  RESUME STARTS HERE  %%%%%%%%%%%%%%%%%%%%%%%%%%%%


\begin{document}

%----------HEADING----------
\begin{center}
    \textbf{\Huge \scshape Prayush B Menon} \\ \vspace{1pt}
    \small Phone Number: 628-263-2031 $|$ Email: \href{mailto:bmenonprayush@gmail.com}{bmenonprayush@gmail.com} \\
    LinkedIn: \href{https://linkedin.com/in/prayush-menon}{linkedin.com/in/prayush-menon} $|$
    GitHub: \href{https://github.com/k0-R0}{github.com/k0-R0} $|$
    Portfolio: \href{https://k0-r0.github.io/}{k0-r0.github.io}
\end{center}

%-----------CAREER OBJECTIVE------------
\section{CAREER OBJECTIVE}
Embedded Systems \& Firmware Engineer with an ECE background and hands-on expertise in bare-metal C/C++ development on PIC microcontrollers. Strong foundation in hardware protocols (CAN, UART, SPI, I2C), memory-efficient data structures, Linux systems programming, and debugging with GDB. Experienced in end-to-end hardware-software integration, multi-threaded POSIX environments, and low-level firmware engineering.

%-----------TECHNICAL SKILLS-----------
\section{TECHNICAL SKILLS}
 \begin{itemize}[leftmargin=0.15in, label={}]
    \item{
     \textbf{Programming Languages}{: C, C++, Embedded C, Java} \\
     \textbf{Microcontrollers \& Protocols}{: PIC18F4580, UART, I2C, SPI, CAN Protocol} \\
     \textbf{Developer Tools \& OS}{: Linux (Fedora), GDB, Make, Git, Neovim, Vim, Cursor} \\
     \textbf{Software Testing \& Management}{: Test Automation, JIRA, Octane}
    }
 \end{itemize}

%-----------PROJECTS AT EMERTXE-----------
\section{PROJECTS AT EMERTXE}
  \resumeSubHeadingListStart

    \resumeSubheading
      {Minishell (Unix Command Shell)}{\href{https://github.com/k0-R0/Minishell}{Github}}
      {C, Linux Systems Programming, IPC, Process Management, Signals}{}
      \resumeItemListStart
        \resumeItem{\textbf{Project brief:} Built a custom Unix command shell in C to understand how Linux manages processes and system calls. The shell supports executing programs, internal built-in commands, multi-stage pipes, signal handling, custom prompts, and job control.}
        \resumeItem{\textbf{Key challenges \& Learnings:}}
        \resumeItemListStart
            \resumeItem{Chained multiple commands using \texttt{pipe()} and redirected standard streams with \texttt{dup2()}, synchronizing process completion using \texttt{waitpid()}.}
            \resumeItem{Added signal handling for \texttt{SIGINT}, \texttt{SIGTSTP}, and \texttt{SIGCHLD} to suspend foreground jobs and resume them in the background or foreground (\texttt{bg}/\texttt{fg}).}
            \resumeItem{Used a singly linked list to track process IDs and execution states for background and stopped jobs.}
        \resumeItemListEnd
      \resumeItemListEnd

    \resumeSubheading
      {Multithreaded TFTP Server \& Client}{\href{https://github.com/k0-R0/TFTP}{Github}}
      {C, UDP Sockets, POSIX Threads, Semaphores, Systems Programming}{}
      \resumeItemListStart
        \resumeItem{\textbf{Project brief:} Built a multithreaded TFTP file transfer server and client in C over UDP sockets. The server handles multiple file uploads and downloads at the same time by spawning worker threads, using a counting semaphore to prevent server overload and a stop-and-wait protocol for reliable delivery over UDP.}
        \resumeItem{\textbf{Key challenges \& Learnings:}}
        \resumeItemListStart
            \resumeItem{Capped active worker threads at 10 using a counting semaphore (\texttt{sem\_t}); used non-blocking \texttt{sem\_trywait()} so the main server loop immediately replies with a ``Server is busy'' packet when full without freezing.}
            \resumeItem{Gave each worker thread its own ephemeral UDP socket to avoid packet collisions, and configured threads as detached (\texttt{PTHREAD\_CREATE\_DETACHED}) so memory and resources are freed automatically on exit.}
            \resumeItem{Built a stop-and-wait acknowledgment protocol with custom packet structs (\texttt{cmd\_packet}, \texttt{data\_packet}, \texttt{ack\_packet}) to track 512-byte block sequence numbers and resend unacknowledged blocks over UDP.}
            \resumeItem{Supported 3 transfer modes---standard binary (\texttt{MODE\_OCTET}), single-byte blocks, and a mail mode that converts \texttt{\textbackslash n} to \texttt{\textbackslash n\textbackslash r}---along with an interactive CLI client for batch file transfers.}
        \resumeItemListEnd
      \resumeItemListEnd

  \newpage

    \resumeSubheading
      {Car Black Box (Event Data Recorder)}{\href{https://github.com/k0-R0/car_blackbox}{Github}}
      {Embedded C, PIC18F4580, I2C, UART, ADC, EEPROM, RTC, LCD}{}
      \resumeItemListStart
        \resumeItem{\textbf{Project brief:} Built an event data recorder on a PIC18F4580 microcontroller in Embedded C to log vehicle telemetry during driving and crash events. The system records timestamps from an external RTC, gear shifts, and speed potentiometer values, saving the 10 most recent records in an external EEPROM with an LCD menu and serial PC export.}
        \resumeItem{\textbf{Key challenges \& Learnings:}}
        \resumeItemListStart
            \resumeItem{Wrote bare-metal I2C master routines to communicate with the DS1307 real-time clock and AT24C04 EEPROM on a shared two-wire bus.}
            \resumeItem{Built a non-volatile circular buffer in EEPROM that retains the 10 latest event records across power resets.}
            \resumeItem{Wrote drivers for ADC potentiometer speed sampling, matrix keypad debouncing, a 16x2 LCD menu, and UART serial communication at 9600 baud.}
        \resumeItemListEnd
      \resumeItemListEnd

    \resumeSubheading
      {Image Steganography (LSB)}{\href{https://github.com/k0-R0/Steganography}{Github}}
      {C, Bit Manipulation, Image Processing, Binary File I/O}{}
      \resumeItemListStart
        \resumeItem{\textbf{Project brief:} Built a command-line tool in C to hide and extract text messages inside 24-bit BMP images using the Least Significant Bit (LSB) technique. The program verifies the BMP header and embeds message bits into RGB pixel bytes without visibly changing the image.}
        \resumeItem{\textbf{Key challenges \& Learnings:}}
        \resumeItemListStart
            \resumeItem{Used bit masking and shifting to replace the least significant bits of image pixel data with message payload bits.}
            \resumeItem{Protected the 54-byte BMP header to keep file metadata intact and ensure the modified image opens in standard image viewers.}
            \resumeItem{Stored message size and magic string signatures inside the image so the decoder can verify and extract the exact secret data.}
        \resumeItemListEnd
      \resumeItemListEnd

    \resumeSubheading
      {Car Dashboard (Automotive CAN Simulator)}{\href{https://github.com/k0-R0/Car_Dashboard}{Github}}
      {C, Embedded Systems, PIC18F4580, CAN Protocol, LCD}{}
      \resumeItemListStart
        \resumeItem{\textbf{Project brief:} Built an automotive dashboard simulator using two PIC18F4580 microcontrollers communicating over a CAN bus. One microcontroller acts as an engine ECU broadcasting speed, RPM, and gear status, while the second acts as a dashboard ECU and renders the data on an LCD.}
        \resumeItem{\textbf{Key challenges \& Learnings:}}
        \resumeItemListStart
            \resumeItem{Configured PIC18F4580 CAN registers, message identifiers, and interrupts to broadcast vehicle telemetry frames.}
            \resumeItem{Wrote display logic on the receiver ECU to filter incoming CAN frames, unpack data bytes, and update the dashboard LCD in real time.}
            \resumeItem{Simulated driving state changes and sensor inputs to test synchronized communication across the CAN bus.}
        \resumeItemListEnd
      \resumeItemListEnd

    \resumeSubheading
      {Lexical Analyzer for C}{\href{https://github.com/k0-R0/Lexical-Analyzer}{Github}}
      {C, Lexical Analysis, Tokenization, Binary Search, File I/O}{}
      \resumeItemListStart
        \resumeItem{\textbf{Project brief:} Built a C source code tokenizer that reads files character-by-character, strips comments and preprocessor lines, and groups code into keywords, identifiers, constants, strings, and operators.}
        \resumeItem{\textbf{Key challenges \& Learnings:}}
        \resumeItemListStart
            \resumeItem{Used binary search ($O(\log N)$) over sorted symbol arrays to quickly look up and identify C keywords and operators.}
            \resumeItem{Built parser state logic to cleanly remove single-line (\texttt{//}) and multi-line (\texttt{/* */}) comments while preserving line numbers.}
            \resumeItem{Added syntax validation to catch unmatched parentheses, brackets, and braces (\texttt{()}, \texttt{\{\}}, \texttt{[]}), as well as unclosed strings.}
        \resumeItemListEnd
      \resumeItemListEnd

  \newpage

    \resumeSubheading
      {Inverted Search Engine}{\href{https://github.com/k0-R0/Inverted_SearchEngine}{Github}}
      {C, Hash Tables, File I/O, Data Structures}{}
      \resumeItemListStart
        \resumeItem{\textbf{Project brief:} Built a text indexing and search engine in C. The program scans text files, builds an inverted index using a custom hash table, and maps each word to a linked list of file occurrences for fast keyword lookups.}
        \resumeItem{\textbf{Key challenges \& Learnings:}}
        \resumeItemListStart
            \resumeItem{Implemented a hash table with linked-list collision handling to index and retrieve search terms in $O(1)$ average time.}
            \resumeItem{Parsed multiple text files to track word counts, file names, and sentence positions without duplicating memory.}
            \resumeItem{Added file persistence routines to save the generated database to disk and reload it without re-parsing files.}
        \resumeItemListEnd
      \resumeItemListEnd

    \resumeSubheading
      {Source2HTML (C Source to HTML Syntax Highlighter)}{\href{https://github.com/k0-R0/Source-2-HTML}{Github}}
      {C, State Machines, Lexical Analysis, File I/O, HTML/CSS}{}
      \resumeItemListStart
        \resumeItem{\textbf{Project brief:} Built a C program that converts C source code into syntax-highlighted HTML files. It scans code using a finite state machine, classifies tokens (keywords, comments, strings, preprocessor lines), and wraps them in CSS-styled HTML tags.}
        \resumeItem{\textbf{Key challenges \& Learnings:}}
        \resumeItemListStart
            \resumeItem{Built a state machine with distinct states (idle, string, comment, directive, number) to accurately parse complex C syntax.}
            \resumeItem{Mapped tokens to CSS color classes to highlight data types, control keywords, and preprocessor directives distinctly.}
            \resumeItem{Escaped special HTML characters (\texttt{\&lt;}, \texttt{\&gt;}, \texttt{\&amp;}) to prevent browser rendering errors, and added an option for line numbering.}
        \resumeItemListEnd
      \resumeItemListEnd

    \resumeSubheading
      {MP3 Tag Reader \& Editor}{\href{https://github.com/k0-R0/MP3_TagReader}{Github}}
      {C, Binary Parsing, File I/O, ID3v2 Tags}{}
      \resumeItemListStart
        \resumeItem{\textbf{Project brief:} Built a command-line tool in C to inspect and edit ID3v2 metadata (title, artist, album) inside MP3 files without altering the audio data.}
        \resumeItem{\textbf{Key challenges \& Learnings:}}
        \resumeItemListStart
            \resumeItem{Parsed ID3v2.3 binary frame headers and calculated variable-length frame sizes to navigate file offsets.}
            \resumeItem{Handled big-endian byte conversions to accurately read frame size descriptors from binary files.}
            \resumeItem{Safely updated tag metadata in-place within the file header without corrupting the underlying MP3 audio stream.}
        \resumeItemListEnd
      \resumeItemListEnd

    \resumeSubheading
      {Arbitrary Precision Calculator (APC)}{\href{https://github.com/k0-R0/APC}{Github}}
      {C, Data Structures, Algorithms, Dynamic Memory Allocation}{}
      \resumeItemListStart
        \resumeItem{\textbf{Project brief:} Built an arbitrary-precision calculator in C using doubly linked lists, allowing mathematical operations on numbers of any size without being limited by 64-bit data types.}
        \resumeItem{\textbf{Key challenges \& Learnings:}}
        \resumeItemListStart
            \resumeItem{Wrote arithmetic algorithms for addition, subtraction, multiplication, and division across linked list nodes, handling carries and borrows digit-by-digit.}
            \resumeItem{Handled dynamic memory allocation and cleanup carefully to ensure zero memory leaks during large arithmetic calculations.}
            \resumeItem{Handled edge cases including sign management (positive/negative results), zero padding, and division-by-zero checks.}
        \resumeItemListEnd
      \resumeItemListEnd

    \resumeSubheading
      {Address Book}{\href{https://github.com/k0-R0/AddressBook}{Github}}
      {C, Data Structures, File I/O, CLI}{}
      \resumeItemListStart
        \resumeItem{\textbf{Project brief:} Built a command-line contact management tool in C to store, search, edit, and delete contacts, saving all entries to a CSV file for persistent storage.}
        \resumeItem{\textbf{Key challenges \& Learnings:}}
        \resumeItemListStart
            \resumeItem{Organized contacts in dynamically allocated arrays with search and sort functions by name, phone, and email.}
            \resumeItem{Handled user input safely with string bounds checking to prevent buffer overflows and invalid entries.}
            \resumeItem{Built file I/O routines to save contact data in a structured CSV format and reload it when the application starts.}
        \resumeItemListEnd
      \resumeItemListEnd

  \resumeSubHeadingListEnd

  \newpage

%-----------WORK EXPERIENCE-----------
\section{WORK EXPERIENCE}
  \resumeSubHeadingListStart

    \resumeSubheading
      {Advanced Embedded Systems Course}{July 2024 -- Ongoing}
      {Emertxe Information Technologies}{Bangalore, India}
      \resumeItemListStart
        \resumeItem{Undergoing comprehensive technical training program aligned with Skill India / NSDC.}
      \resumeItemListEnd

    \resumeSubheading
      {Senior Quality Engineer I}{July 2024 -- January 2026}
      {Fresh Prints}{Remote}
      \resumeItemListStart
        \resumeItem{Built a frontend performance testing infrastructure using Lighthouse, reducing manual testing time by 80\% through automated bi-weekly performance monitoring across all environments.}
        \resumeItem{Led data migration validation for 7000+ design records from Shopify to an actual CMS - Directus, developing custom validation scripts that ensured 100\% data integrity during platform transition.}
        \resumeItem{Identified critical P1 security vulnerabilities through pen testing tools like Burpsuite that allowed unauthorized client data access, collaborating with the development team for immediate resolution.}
      \resumeItemListEnd

    \resumeSubheading
      {QA Automation Engineer}{June 2022 -- July 2024}
      {Cognizant Technology Solutions}{Bangalore, India}
      \resumeItemListStart
        \resumeItem{Engineered automated test suites for web applications, scaling test automation coverage and validating core business logic across production releases.}
        \resumeItem{Automated end-to-end user journeys across multi-tier applications, improving test execution efficiency and platform reliability.}
      \resumeItemListEnd

  \resumeSubHeadingListEnd

%-----------COURSE WORK-----------
\section{COURSE WORK}
  \resumeSubHeadingListStart
    \resumeItem{Microprocessor, Digital Electronics, Digital Signal Processing, Embedded Systems Design}
  \resumeSubHeadingListEnd

%-----------EDUCATION-----------
\section{EDUCATION}
  \resumeSubHeadingListStart
    \resumeSubheading
      {Amrita School of Engineering}{Kollam, Kerala}
      {B.Tech in Electronics and Communications Engineering (CGPA: 8.1/10)}{June 2018 -- August 2022}
  \resumeSubHeadingListEnd

\end{document}
